\section{Why metagenomic read classification challenges current GFMs}
\label{sec:challenges}

Three properties of the task make it a poor fit for foundation models designed
and evaluated on other genomic problems.

\paragraph{Short reads truncate long-range context.}
Modern GFMs derive much of their advantage from modelling long-range
dependencies---HyenaDNA and Evo operate over tens to hundreds of kilobases
\citep{Nguyen2023hyena,Nguyen2024,Brixi2025}. A 150\,bp Illumina read offers
none of that context. Whatever discriminative information exists must be
extracted from local composition, so the mechanism that distinguishes a strong
model from a weak one on genome-scale tasks is largely unavailable here. This
predicts that input representation (how 150\,bp is turned into tokens) should
matter more than backbone depth or pre-training corpus---a prediction we test
directly.

\paragraph{Tokenization discards or preserves local signal.}
The taxonomic signal in a short read lives in its $k$-mer composition. How a
model tokenizes that read therefore determines how much signal reaches the
network. Nucleotide Transformer v2 uses \emph{non-overlapping} 6-mers: a
150\,bp read becomes $\sim$25 tokens, and any motif that straddles a 6-mer
boundary is split across two tokens whose identity depends on frame. Kraken\,2,
by contrast, effectively uses long \emph{overlapping} $k$-mers ($k\approx 35$),
capturing every length-$k$ substring. Overlapping long-$k$ tokenization is much
closer to the representation that classical metagenomics has always relied on,
but it is the exception rather than the rule among pre-trained GFMs. This gap
between ``how GFMs tokenize'' and ``what the task needs'' is, we argue, the
central issue, and it motivates treating tokenization as a first-class
experimental variable.

\paragraph{Large, long-tailed label spaces.}
Real communities span hundreds of genera and thousands of species with
abundances covering several orders of magnitude. A read-level classifier must
therefore solve a 120-way (genus) or 1{,}535-way (species) problem from 150\,bp,
with training data that is intrinsically imbalanced. Naively rebalancing the
classes trades per-read accuracy and real-abundance fidelity for macro-averaged
metrics (Section~\ref{sec:scaling}), so the choice of evaluation metric is not
cosmetic---it changes which model looks best. This is why we evaluate at two
levels: per-read accuracy, which measures raw discriminative power, and
sample-level abundance, which measures the quantity practitioners actually
report.

\medskip
\noindent Together these properties imply that a fair evaluation must (i) hold
the architecture fixed while varying tokenization, (ii) separate data scale
from model capacity, and (iii) report both read-level and sample-level
outcomes. The benchmark in Section~\ref{sec:design} is built around exactly
these requirements.
